\section{Further research questions and a proposed programme}
\label{sec:research}

The next stage should be driven by the parameters that remain in the
proved bounds. The following questions are proposed research tasks,
not assertions that their desired answers hold. A useful negative answer
would often be as valuable as a positive one, because it would prevent
investment in a compression invariant that cannot control the algorithm.

\subsection{Priority I: the global structural obstruction}

\begin{question}[Connected components after exact reductions]
\label{q:components}
For which diagram families, or for which residual inputs after the
existing certified filters, can one prove a bound on the largest reduced
connected component size $b$ that makes
$b^2 B_m=O((\log n)^2)$?
\end{question}
This is the most direct missing theorem for the report-07 route. At fixed
frontier, the sufficient condition permits $b=O(\log n)$; it does not
permit arbitrary polynomial growth in $b$. Experiments should record
component sizes, homological spans, matching multiplicities, dot-degree
layers, and weights at every stage. Families designed to resist
connected-sum factorization and elementary simplification are especially
valuable. If $b$ grows too quickly even at fixed width, determine whether
the growth is forced under every admissible cancellation order, or is
an artifact of the chosen pivot policy. A proof of such an obstruction
would redirect the programme toward stronger equivalence relations.

\begin{question}[Compression beyond identical serialized components]
Can a constructive chain-homotopy normalization replace large connected
components by a bounded description while preserving the recognition
criterion and efficient future crossing updates?
\end{question}
Exact component sharing only merges identical keys. It does not perform
arbitrary changes of basis or discover every homotopy equivalence.
Candidate descriptions include graded module decompositions, interval
structures in restricted categories, and presentations over small
endomorphism algebras. Each proposal must specify the allowed basis
changes, provide an algorithm to compute the description, and bound
composition after a crossing is added. Merely knowing that equivalent
complexes have a small representative is insufficient if finding or
updating it requires enumerating exponentially many bases. Test the
proposal first on the known multiplicity families and on connected
examples that cannot be removed by visible sum decomposition.

\begin{question}[Crossing order optimized for the real cost]
Can crossing orders be constructed with controlled
$w+\log(N+2)$, or controlled $b^2 B_{w/2}$, rather than with small
frontier alone?
\end{question}
A frontier objective sees only one part of the resource envelope.
The stress fixture in Table~\ref{tab:scanner} has a small frontier and
thousands of surviving objects. A useful ordering certificate would
explain why its chosen local pieces have bounded component growth, not
just count exposed endpoints. Investigate separator decompositions with
independent frontier components, preserving a certified disk order only
where that order actually exists. Any Catalan estimate must carry the
topological ordering certificate that makes it valid. Approximation or
heuristic ordering results should charge the time spent finding the order.

\begin{question}[A completeness theorem for a portfolio]
Can one prove that every diagram either admits a cheap structural or
finite-target certificate, or enters the homology backend with parameters
satisfying a quasi-polynomial bound?
\end{question}
This formulation respects the difference between complete homology
calculation and unknot recognition. The full invariant can be expensive
even when a simple certificate already decides the knot. A successful
theorem would describe a finite or effectively generated set of cases,
prove coverage, and bound how the appropriate case is found. It must also
account for unknot inputs, where knottedness filters cannot help. The
current three-braid, homogeneous-diagram, connected-sum, Alexander,
Jones, and finite-target stages are concrete candidates for parts of
such a partition, but present experiments do not prove that their
residual class has small $b$ or small width.

\subsection{Priority II: exploit the grading in the representation}

\begin{question}[Rank-layer coefficient storage]
Can a homogeneous $k$-circle coefficient of dot degree $d$ be stored
and manipulated in time near $\binom{k}{d}$, with no full $2^k$ table
allocated during contraction or output reconstruction?
\end{question}
Theorem~\ref{thm:homogeneous} reduces the transform arithmetic, while
Theorem~\ref{thm:types} reduces the number of possible coefficients.
The current implementation still allocates subset-indexed tables.
Combinatorial ranking of fixed-cardinality subsets, restricted zeta
transforms, or sparse-to-layer conversions might realize the counting
gain in memory as well. The challenge is closure: identifying input
variables can kill terms and combine them by parity, and the output map
has nontrivial block structure. A valid result must cover those maps,
not just multiplication of two already indexed middle layers. Compare
the resulting cost with both input and output layer sizes.

\begin{question}[A grading-aware scanner interface]
Should the quantum shifts be stored explicitly to eliminate impossible
entries earlier, restrict pivot search, and specialize every genuine
cancellable unit to $1$?
\end{question}
The proof shows that the information is recoverable; the current audit
reconstructs it separately. An explicit grading can reject incompatible
matrix entries before compiling a surface plan and can provide a useful
internal consistency check. Measure whether the saved work exceeds the
cost of maintaining shifts and layer metadata. The interface must retain
the generic algebra fallback for clients that intentionally supply
inhomogeneous morphisms. A separate graded object type is preferable to
silently assuming homogeneity from the packed coefficient's shape.

\begin{question}[Remove sparse-case adapter overhead]
Can the implementation preserve the baseline cost on small or sparse
entries while retaining the proved dense envelope?
\end{question}
This is a high-confidence engineering task suggested directly by the
paired timings. A fixed small-frontier dispatch to the original primitive
preserves the asymptotic theorem, because its cutoff is constant.
More general dispatch must also charge output expansion and large-integer
updates; a support-only estimate can be misleading. Profile plan lookup,
support extraction, homogeneity detection, byte conversion, and cache
reuse separately. Promote a new default only after paired scans on both
ordinary inputs and genuinely dense graded stages demonstrate a useful
tradeoff. Retain raw distributions rather than optimizing one selected
median on a noisy machine.

\begin{question}[What dense entries can actually occur?]
Is there an explicit knot-diagram family producing large, dense
homogeneous differential entries after the prescribed cancellations,
and can their density be predicted from topology?
\end{question}
The arbitrary-algebra family in Section~\ref{sec:convolution} proves a
local worst-case improvement. It does not supply a knot family realizing
that local worst case. A useful construction would preserve exact stage
data and provide a proof of its coefficient support, rather than merely
an opaque slow example. Conversely, a theorem limiting support density
in all relevant surviving complexes might make a different sparse
algorithm the right asymptotic target. This question connects the
coefficient theorem to end-to-end usefulness.

\begin{question}[Other characteristics and integral coefficients]
How far does the plan factorization extend with controlled bit complexity
over odd finite fields or over the integers?
\end{question}
Ranked subset convolution works over much more general coefficient rings
than $\F$~\cite{bhkk}. The geometric specialization needs separate
attention: the handle multiplier $2x$ no longer vanishes in odd
characteristic, integer signs must be consistent, and integer coefficient
growth must be bounded. A graded unit may be a nonzero scalar requiring
inversion, rather than the unique scalar $1$. An extension should specify
the exact Frobenius algebra, prove the new component evaluation formula,
and charge coefficient bit lengths. It should not simply reuse the
characteristic-two positive-genus shortcut.

\subsection{Priority III: bounded witnesses and target selection}

\begin{question}[Constructive seed-plan quality]
Can the greedy $O(n^3)$ seed-plan construction be replaced by a faster
dependency-closure algorithm, with a useful approximation or structural
bound on the resulting seed count?
\end{question}
For fixed targets, lowering $b$ by one saves an exponential factor in
the search parameter, while a smaller planning constant matters for
sub-millisecond fixtures. Maintain the crossing dependencies exposed by
each newly colored arc, and reuse closure information between candidate
seeds. Prove the setup cost of this reuse and keep the final plan
independently checkable. The equality between the minimum Wirtinger
number over diagrams and bridge number~\cite{blair} is a structural
guide; it does not make minimum-seed search on an arbitrary presented
diagram free.

\begin{question}[Adaptive finite targets with a completeness bound]
Can target groups or quandles be constructed from a diagram so that
every nontrivial residual knot has a witness with
$b\log(q+1)=O((\log n)^2)$, and both target construction and witness
search fit the same budget?
\end{question}
The palette obstruction rules out a universal bounded palette. The
NP-completeness result for fixed-target coloring~\cite{kuperbergsamperton}
also prevents a casual inference from a small target to an easy search
for every diagram. A constructive target-selection theorem would need
to exploit special structure or permit a different obstruction when
a target search is difficult. Count the target description, conjugation
table or on-demand operation cost, seed plan, and proof of nonconstancy.
Do not replace all these requirements by an existence theorem for some
finite nonabelian quotient.

\begin{question}[Sharp detecting-quandle size for determinant-one torus knots]
What is the smallest size $q(p)$ of a finite quandle that nontrivially
colors $T(3,p)$, for prime $p>3$? Can the lower bound $q(p)\geq p$
be improved, and can one give explicit detecting targets with comparable
upper bounds?
\end{question}
The proof uses the exponent of the inner group and works for nonfaithful
and disconnected quandles. It does not show that a size-$p$ target exists.
Study conjugacy classes in explicit finite quotients of the torus group,
and compare their class size with the size of the entire group. Distinguish
the number of quandle elements from the encoding length of a succinct
group. Determinant one already removes familiar affine-coloring routes;
the non-affine constructions discussed by Fish, Lisitsa, and
Stanovsk\'y~\cite{fish} give relevant precedent. Small computational
classifications could suggest a sharper theorem, but should not be
advertised as a bound for untested primes.

\subsection{Priority IV: alternate routes to a global theorem}

\begin{question}[Detecting slices with a completeness guarantee]
For which broad knot families does nontrivial reduced Khovanov homology
necessarily appear within a provably shallow extremal homological band?
\end{question}
Combine a genuine detection theorem with the prescribed nice-scan
construction of Kelom\"aki and Sch\"utz~\cite{kelomaki}. For the
displayed slice envelope, $g=O((\log n)^2)$ and $k=O(\log n)$ would
fit the stronger $n^{O(\log n)}$ scale, provided the scan is constructed
within that cost. A fast slice without a theorem locating the required
homology is only a partial obstruction. Also quantify what happens when
the useful grading range is specified in normalized knot-theoretic
degrees rather than unshifted cube degrees; crossing-sign shifts must
be handled explicitly.

\begin{question}[A complete hierarchy cost ledger]
Can the hierarchy approach be organized into explicit compressed data
structures for which every visit, cut, normalization, and violating-disc
step has a proved bit bound?
\end{question}
Lackenby's bound $L(g+1)^L$ has logarithm
$O(\log(L+1)+L\log(g+1))$~\cite{lackenby2026}. Thus
$L\log(g+1)=O((\log n)^2)$ is a natural target for the stronger
scale in this project, while $O((\log n)^3)$ would still yield the
broader quasi-polynomial scale announced in the 2021 handout.
If $g$ is polynomial in $n$, these correspond to hierarchy lengths of
order $\log n$ and $(\log n)^2$, respectively. These statements bound
visit counts only. Compressed surface coordinates can encode extremely
large geometric objects, so naive expansion during a cut would destroy
the argument. The normal-curve algorithms in~\cite{lackenbycurves}
suggest useful primitives, but the full ledger and construction theorem
must be written and proved.

\begin{question}[Certificates for the algebraic reduction itself]
Can every crossing extension and cancellation emit a compact reduction
certificate whose verifier is substantially simpler than the optimized
scanner?
\end{question}
The finite-target work demonstrates the practical value of a search and
verifier separation. For homological reduction, record the pivot,
inclusion and projection maps, and the relevant chain-homotopy identity.
The challenge is certificate size: an explicit transcript can be as
large as the computation. Investigate compressed transcripts with
deterministically replayable local claims, and state whether verification
requires the same coefficient tables. A formalization can begin with the
squarefree-algebra factorization, ranked convolution, and scalar
cancellation identity, while clearly treating the topological
unknot-detection theorem as a separately cited assumption.

\subsection{Concrete milestones}

The next release can have measurable goals without pretending to settle
the global problem immediately.
\begin{enumerate}
\item \textbf{Instrument the surviving structure.} Merge grading metadata
with report-07 sharing in an isolated branch, recording $m$, $b$, $N$,
component weights, and dot-degree distributions. First reproduce the
existing exact outputs; then publish the component-growth data.
\item \textbf{Improve ordinary workloads.} Add a constant-frontier legacy
dispatch and a measured cache policy, while preserving the proved
coefficient envelope. Require paired ordinary scans and genuine graded
dense cases before changing defaults.
\item \textbf{Prove one nontrivial structural family theorem.} Choose a
family after the existing certified filters and prove its actual
component-size or homotopy-description bound. A theorem with an explicit
construction is more informative than a larger table of favorable
timings.
\item \textbf{Publish a falsifiable global route.} State exactly which
structural dichotomy or hierarchy construction would imply the desired
bound, and keep every unproved lemma visible. Search deliberately for
counterexamples to those lemmas before expanding the implementation.
\end{enumerate}

The present continuation has made coefficient operations essentially
linear in their explicit table size up to polynomial frontier factors,
and it has improved the available counting argument through a recovered
grading. The proposed programme now concentrates on the amount and
structure of data that the algorithm must construct. That is the point
at which the quasi-polynomial objective will be decided.
